%=============================================================================!
% 11.6  HEAVIER HYDROGENS                                                     !
%=============================================================================!
\section{Heavier hydrogens}
\label{sec:cs-mass}

The potential energy depends only on where the atoms are; the masses
enter only through Newton's law, and set how fast the
atoms respond to the forces. A third way to lengthen the step changes the
masses so that the fastest motions slow down, while keeping the energy
surface, and with it the forces, as it was. This section works out what
the change does to the vibrations of water and measures the step it buys.

\subsection*{Moving mass}

A bond between two atoms vibrates with $\omega = \sqrt{k/m_{\mathrm r}}$
(\cref{sec:os-bodies}), and for a bond to hydrogen the reduced mass is
close to the hydrogen's own: for O--H, $m_{\mathrm r} = 15.999\times1.008/
(15.999 + 1.008) = \SI{0.9483}{\amu}$. Hydrogen mass repartitioning makes
each hydrogen heavier and takes the same mass from the atom it is bonded
to, so that every molecule keeps its total mass. Tripling the hydrogen's
mass to \SI{3.024}{\amu} takes $2\times2.016$ from the oxygen of water,
leaving \SI{11.967}{\amu}; the molecule still weighs \SI{18.015}{\amu}, so
its centre of mass, now a little nearer the hydrogens
(\cref{ex:cs-hmr-centre}), has the same acceleration, the total force
over the total mass, as before. The O--H reduced mass becomes
$11.967\times3.024/(11.967 + 3.024) = \SI{2.4140}{\amu}$, 2.55 times
larger, and the stretch slows by $\sqrt{2.55} = 1.60$. The function
\texttt{repartition\_masses} of \texttt{mdlab.constraints} makes the
change for any list of bonds to hydrogen.

The normal modes of the spring model of water (\cref{sec:ro-freedom})
with the new masses show a similar slowing, by 1.68 for the fastest
vibration, which falls from
$\omega = 0.857$ to \SI{0.509}{\radian\per\femto\second}, its period rises
from 7.33 to \SI{12.35}{\femto\second}, and the stability limit $2/\omega$
of \cref{sec:in-stability} from 2.33 to \SI{3.93}{\femto\second}; the bend
slows from a period of 20.9 to \SI{32.2}{\femto\second}, and with the O--H
bonds held from 17.0 to \SI{27.2}{\femto\second}. A rigid molecule
cannot vibrate, but it turns, and moving mass outwards from the oxygen to
the hydrogens raises its moments of inertia (\cref{sec:ro-tensor}): the
principal moments, at the TIP3P shape, go from 0.615, 1.155 and \SI{1.770}{\amu\,\angstrom
\squared} to 1.379, 3.465 and \SI{4.844}{\amu\,\angstrom\squared}, so its
turning in the cage of its neighbours slows too.

\subsection*{What it buys}

The open markers of \cref{fig:cs-step}(a) measure it on the box of 64
molecules. At the 1\% level, water with its O--H bonds held goes from
0.97 to \SI{1.31}{\femto\second} with heavier hydrogens, a factor of
1.35, although its bend, the fastest motion left, slows by 1.60, so the
bend is not all that sets its step; rigid water goes from 1.98 to
\SI{3.38}{\femto\second}, a factor of 1.71, and is then 8.2 times the
step of the flexible model with ordinary masses. Hopkins and co-workers,
citing earlier studies in which repartitioning let the step be increased
by up to a factor of two, ran a short chain of three amino acids and a
protein with repartitioned masses and steps of
\SI{4}{\femto\second}\cite{Hopkins2015}.

\subsection*{What it keeps, and what it changes}

Repartitioning leaves the potential energy and the forces unchanged, and
Chapter~12 shows that the probability of finding the atoms in a given
arrangement, at a given temperature, depends on the potential energy
alone and not on the masses. Averages over the arrangements of the atoms,
such as their mean potential energy or how the molecules are packed, are
therefore unchanged, exactly so for molecules that are free or held wholly
rigid; with only some distances held, as with the O--H bonds alone, the
held distances make the masses enter the probability of the angles left
free, a small effect that this book does not pursue. Everything that depends on how fast the atoms move is
changed: heavier hydrogens move more slowly, so the rates at which
molecules turn and diffuse, and the frequencies of their vibrations,
belong to a system with different masses. Hopkins and co-workers report no
significant difference in the rates and the equilibrium properties they
computed with \SI{4}{\femto\second} steps and repartitioned
masses\cite{Hopkins2015};
that is a finding about their systems and properties, and a study of
motion at the atomic scale should check its own.

\begin{keyidea}
Moving mass from heavy atoms to their hydrogens keeps the total mass and
the forces, and raises the reduced masses of the fastest bonds: tripling
hydrogen's mass in water slows the O--H stretch by a factor of 1.60. On
64 molecules it lengthens the step at 1\% by 1.35 with only the O--H bonds
held and by 1.71 for rigid water. Averages over arrangements are kept,
exactly for free or rigid molecules; rates of motion are not.
\end{keyidea}

\notebook{11}{choose the hydrogen mass and watch the vibrations slow}

\begin{selfcheck}
In a C--H bond, carbon \SI{12.011}{\amu} and hydrogen \SI{1.008}{\amu}, the
hydrogen is made three times heavier at the carbon's expense. By what
factor does the stretch slow?
\end{selfcheck}
